\documentclass[10pt,a4paper]{amsart}
\usepackage[margin=19mm]{geometry}
\usepackage{amsmath,amssymb,mathtools}
\usepackage[scr=rsfs,cal=euler]{mathalfa}
\usepackage{xfrac}
\usepackage[unicode,hidelinks,bookmarks=false]{hyperref}
\newcommand{\ie}{i.e.\ }
\newcommand{\cf}{cf.\ }
\newcommand{\resp}{respectively\ }
\newcommand{\Subsection}{\subsection}
\providecommand{\nobreakdash}{\nobreak\hbox{-}}
\setlength{\emergencystretch}{2em}
\multlinegap0pt
\usepackage{amssymb}
\usepackage{dsfont}
\usepackage{algorithm}\usepackage{algpseudocode}
\usepackage{mathtools}
\usepackage[shortlabels]{enumitem}
\newtheorem{lemma}{Lemma}
\newtheorem{proposition}{Proposition}
\title{Pointwise Convergence Analysis for Approximations of Optimal Transport Problems with a Target Measure that Has Unbounded Support}
\author{Axel G. R. Turnquist}
\date{Source: arXiv:2603.01645v1 (CC BY 4.0)}
\begin{document}
\maketitle

\noindent\emph{Source and licence.} Pointwise Convergence Analysis for Approximations of Optimal Transport Problems with a Target Measure that Has Unbounded Support. Version arXiv:2603.01645v1 (CC BY 4.0), distributed by arXiv under the Creative Commons Attribution 4.0 International licence, http://creativecommons.org/licenses/by/4.0/, as recorded in the arXiv licence metadata for this exact version and shown on the abstract page https://arxiv.org/abs/2603.01645v1. Full licence notice: \texttt{licenses/arxiv-2603.01645-CC-BY-4.0.txt}.\par\smallskip

\noindent\textit{Editorial scope.} Counted unit: the article's proposition prop:logconcave (source0.tex lines 853--859). The article's complete original proof of it is delivered with it (source0.tex lines 861--918, one \texttt{proof} environment of 58 lines). No mathematics is written, repaired or re-derived: the statement and the proof are the article's own lines, in the article's own notation, and no retained line is substituted or re-flowed.  Retained uncounted as the source-local context the proof needs: lines 755--760; lines 757--759.  Nothing was omitted from any retained span, so the package carries no omission record.  No retained line contains a citation, so the wrapper carries no bibliography and no bibliography entry.  No retained line contains a figure environment, an image-inclusion command or drawing code, so no image is delivered and the package declares no asset dependency.  Editorial formatting only, in addition: an AMS \texttt{amsart} preamble that restates the article's own declarations and macros used by the retained lines, namely the environments \texttt{lemma}, \texttt{proposition} on separate counters, together with the standard packages those lines need.  The retained environments print in this document as Lemma 1, Proposition 1 in order of appearance, whereas the article prints the same items with the numbering of its own full document.  The complete native source of the article as distributed in the arXiv e-print for this exact version is bundled unchanged as \texttt{sources/195568/source0.tex}.
\begin{lemma}\label{lemma:firstbound}
Let $\nu \in \mathcal{P}(\mathbb{R}^n)$. Then, the following bounds the $W_1$ distance between $\nu$ and $\nu_R$:
\begin{equation}\label{eq:generalbound}
W_1(\nu, \nu_R) \leq \int_{\mathbb{R}^n \setminus B_{R}(0)} \vert x \vert d\nu(x) + R(1-\nu(B_{R}(0))).
\end{equation}
\end{lemma}

\begin{equation}
W_1(\nu, \nu_R) \leq \int_{\mathbb{R}^n \setminus B_{R}(0)} \vert x \vert d\nu(x) + R(1-\nu(B_{R}(0))).
\end{equation}

\begin{proposition}\label{prop:logconcave}
Let $\nu$ be a log-concave distribution, i.e. $f_{\nu}(x) = e^{-\varphi(x)}$ for a convex function $\varphi(x)$ and let $R \geq 1$. Then, there exist constants $C = C(\nu, n)$ and $a>0$ such that:
\begin{equation}
W_1(\nu, \nu_{R}) \leq C(\nu, n) R^n e^{-aR},
\end{equation}
where $C(\nu, n)$ is given in Equation~\eqref{eq:fastrate}.
\end{proposition}

\begin{proof}
In Parts A-D, we prove that there exist constants $a>0$ and $b \in \mathbb{R}$ such that $\varphi(x) \geq a \vert x \vert + b$. We use the fact that $\varphi$ is convex and that $\int_{\mathbb{R}^n} e^{-\varphi(x)}dx = 1$. Most of the difficulty of the proof is contained in showing this (Parts A-D). Once we have this bound, we perform a routine computation to get the Wasserstein-1 distance between $\mu$ and $\mu_{R}$ (Part G).

Part A) We first show that $\varphi(x)$ must be a proper convex function. Suppose $\varphi(x)$ is not proper. This means that i) $\varphi(x) = +\infty$ for all points, in which case $e^{-\varphi(x)} = 0$ for all $x\in \mathbb{R}^n$ which contradicts the fact that $\int_{\mathbb{R}^n} e^{-\varphi(x)}dx = 1$ or ii) $\varphi(x) = -\infty$ for some $x \in \mathbb{R}^n$. Let $E$ denote the set where $\varphi(x) = -\infty$. Then $E$ must be a convex set of Lebesgue measure zero, which is a subset of a hyperplane. Then $\partial E$ is a measurable set since it is the boundary of a convex set, and it has measure zero. The set $\partial E$ is the only set on which $\varphi$ can take finite values. Because of this, then we are forced to conclude that $\int_{\mathbb{R}^n} e^{-\varphi(x)}dx = 0$, which, again, is a contradiction. For a similar reason, $\varphi$ cannot only be finite on a set of measure zero.

Part B) Now that we know $\varphi(x)$ must be proper, let us assume that $\varphi(x)$ has a minimum on $\mathbb{R}^n$ at, say, $x_0 \in \mathbb{R}^n$. We will demonstrate that if the affine dimension of $\text{dom}(\varphi)$ is equal to $n$, then the minima of $\varphi$ must exist on a compact set. Also, by the integrability condition, we cannot have the the affine dimension of $\text{dom}(\varphi)$ be less than $n$.

After an affine coordinate transformation, we assume $\varphi(0) = 0$. If $\varphi(x) = 0$ for some $x \neq 0$, then $\varphi$ must be equal to zero along the line connecting $x$ and $0$, since $\varphi$ is convex. Therefore, if $\varphi(x)$ takes the value zero arbitrarily far away from the origin, it takes the value zero on the ray emanating from the origin. By another affine coordinate transformation, we take this ray to be in the direction of the positive first coordinate axis. Thus, by the choice of our coordinate system, $\varphi(x)=0$ for any $x$ such that $x_1 \geq 0$ and $x_2 = \dots = x_n = 0$. Since in Part A) we proved that $\varphi$ cannot only be finite on a set of measure zero, and $\text{dom}(\varphi)$ is convex, then the affine dimension of $\text{dom}(\varphi)$ must be equal to $n$. Thus, $\text{int}(\text{dom}(\varphi)) \neq \emptyset$. Let $\tilde{x} \in \text{int}(\text{dom}(\varphi))$. Then, there exists a ball $B_{\epsilon}(\tilde{x}) \subset \text{int}(\text{dom}(\varphi))$, and therefore $\varphi(y) <\infty$ for every $y \in B_{\epsilon}(\tilde{x})$. Since $e_1$ is a direction of recession of $\text{dom}(\varphi)$, then $y + te_1 \in \text{dom}(\varphi)$ for $t\geq 0$ and every $y \in B_{\epsilon}(\tilde{x})$. Furthermore, $(e_1, 0) \subset \mathbb{R}^{n+1}$ is a direction of recession of $\text{epi}(\varphi)$ and therefore, $\varphi(y + t e_1) \leq \varphi(y)$ for every $t \geq 0$.  Thus,
\begin{align}
\int_{\mathbb{R}^n} e^{-\varphi(x)}dx &\geq \sum_{k=1}^{\infty} \int_{B_{\epsilon}(\tilde{x} + 2 \epsilon k)} e^{-\varphi(x)}dx \\
&=\sum_{k=1}^{\infty} \int_{B_{\epsilon}(\tilde{x})} e^{-\varphi(x)}dx = \infty.
\end{align}

%Part B) Now that we know $\varphi(x)$ must be proper, let us assume that $\varphi(x)$ has a minimum on $\mathbb{R}^n$ at, say, $x_0 \in \mathbb{R}^n$. After an affine coordinate transformation, we assume $\varphi(0) = 0$. If $\varphi(x) = 0$ for some $x \neq 0$, then $\varphi$ must be equal to zero along the line connecting $x$ and $0$, since $\varphi$ is convex. Therefore, if $\varphi(x)$ takes the value zero arbitrarily far away from the origin, it takes the value zero on the ray emanating from the origin. By another affine coordinate transformation, we take this ray to be in the direction of the positive first coordinate axis. For $x \in \mathbb{R}^n$ with $x_1 \geq 0$, we must have that any supporting tangent plane to the epigraph of $\varphi$ must be of the form $a x + b$, with $a \cdot e_1 \leq 0$. If not, then we could move along the ray and the value of $\varphi(x)$ would have to be arbitrarily large to remain above the tangent planes. Hence, $\varphi(x) \leq \tilde{\varphi}(x)$, where $\tilde{\varphi}(x)$ (not necessarily convex) is constant in its first coordinate for $x_1>0$. But then,
%\begin{align}
%\int_{x_1\geq 0} e^{-\varphi(x)}dx &\geq \int_{x_1\geq0} e^{-\tilde{\varphi}(x)}dx \\
%&= \int_{-\infty}^{\infty} \dots \int_{-\infty}^{\infty} \int_{0}^{\infty} e^{-\tilde{\varphi}(x_1, x_2, \dots, x_n)} dx_1 dx_2 \dots dx_n \\
%&= \infty,
%\end{align}
%for all $x$ such that $x_1>0$. Therefore, $\int_{\mathbb{R}^n} e^{-\varphi(x)}dx = \infty$. This shows that if $\varphi(x)$ has a minimum on $\mathbb{R}^n$, then the minimum will not necessarily be unique, but it must occur in a compact set. This argument can be generalized to show that $\varphi(x)$ cannot be a constant real number along any ray emanating from the minimum.

Part C) Suppose $\varphi(x)$ has no minimum. By an affine change of coordinates, let $\varphi(0) = 0$. Then, $\varphi(x) \rightarrow -\infty$ along a ray $R = \{ay: a\geq 0, ,y \in \mathbb{R}^n\}$ emanating from the origin. Because $\varphi(x)$ is convex, it must be convex along the ray $R$. Since $\varphi(0) = 0$ and $\varphi$ is convex (and therefore monotonically decreasing along the ray $R$), then $\varphi(x) \leq 0$ for $x \in R$.

Then, we define a new function $\tilde{\varphi}(x) = \varphi(x)$ for $x$ such that $\varphi(x) \geq 0$ and $\tilde{\varphi}(x) = 0$ for $x$ such that $\varphi(x) < 0$. Then, $\tilde{\varphi}$ is convex and has a minimum. But, then $\int_{\mathbb{R}^n} e^{-\varphi(x)}dx \geq \int_{\mathbb{R}^n} e^{-\tilde{\varphi}(x)}dx = \infty$, by applying Part B to the function $\tilde{\varphi}(x)$.

Part D) We have thus showed that $\varphi$ must have a minimum. If there do not exist a constant $a>0$ and a vector $b \in \mathbb{R}$ such that $\varphi(x) \geq b + a \vert x \vert$, this must be because $\varphi$ is constant in a ray emanating from the minimum point $x_0$. Without loss of generality, suppose that the minimum of $\varphi$ occurs at the origin. Fix $b  < \varphi(0)$ and $a>0$. Let $S = \{ x: b + a \vert x \vert \leq \varphi(x) \}$. Then, if there is a point $y \in \mathbb{R}^n$ such that $\varphi(y) < b + a \vert y \vert$. Thus, we can find a new value $a'$ that satisfies $a' \leq \frac{\varphi(y) - b}{\vert y \vert}$ and will guarantee that $\forall x \in S$, we have $\varphi(x) \geq b + a' \vert x \vert$. We thus see that the existence of a uniform $a$ that guarantees that $\varphi(x) \geq b + a \vert x \vert$ for all $x \in \mathbb{R}^n$ is equivalent to guaranteeing that $\alpha = \inf_{y} \frac{\varphi(y) - b}{\vert y \vert} > 0$. Since $b<\varphi(x)$ for all $x \in \mathbb{R}^n$, this is only possible if $\varphi(x)$ grows asymptotically slower than any linear function as $\vert y \vert \rightarrow \infty$. The constant $\alpha = 0$ only if $\varphi$ is constant in some ray emanating from the origin (the minimum). This is because if $\varphi$ were not constant along any ray emanating from the origin, then along any ray emanating from the origin, by the convexity of $\varphi$, there would exist a tangent line with nonzero slope lying below $\varphi$ restricted to the ray. Then, since $\varphi$ lies above the tangent line, it grows at least linearly as $\vert y \vert \rightarrow \infty$. To be more precise, fix $z \in \mathbb{S}^{n-1}$ and define $r_{z} = \{ t z : t \geq 0 \}$. Define $\varphi_{z}(t) := \varphi(t z)$. Then, $\varphi_{z}(t)$ is convex on $[0, \infty)$. If $\varphi_{z}$ is not constant on $[0, \infty)$, then since $\varphi_{z}(t)$ is finite for $t = 0$, there exists a value $y >0$ such that $y \in \partial \varphi_{z}(t_0)$ for some $t_0 \in [0, \infty)$. Therefore, $\varphi_{z}(t) \geq \varphi_{z}(t_0) + y(t - t_0)$ for all $t \in [0, \infty)$. Thus, since the value of $\alpha$ can be found by minimizing along rays and then minimizing over $\mathbb{S}^{n-1}$, which is a compact set, the only way for $\alpha$ to equal zero is if $\varphi$ were constant in some ray emanating from the origin. Since above we have proved that this is not possible, we conclude that there must exist $a>0$ and $b \in \mathbb{R}$ such that $\varphi(x) \geq b + a \vert x \vert$ for all $x \in \mathbb{R}^n$.

%Part D) Assume that $\varphi(x)$ is $C^1$ on all of $\mathbb{R}^n$. By an affine coordinate transformation we assume $\varphi(0) = 0$. Now, by the above arguments $\varphi(x)$ takes the value zero in a compact set $M \subset \mathbb{R}^n$. Let $r$ be such that $r > \max_{x \in M} \vert x \vert$. Then, $\nabla \varphi(x) \neq 0$ for $x$ satisfying $\vert x \vert =r$. Set $\alpha = \min_{x \in B_{r}(0)} \vert \nabla \varphi(x) \cdot \hat{x} \vert$. Then, $\alpha>0$, since the minimum is of a continuous function on a compact set. Also, denote $\beta = \max_{x \in B_{r}(0)} \vert \nabla \varphi(x) \cdot \hat{x} \vert$. Then, denote $\rho := \min_{y \in B_{r}(0)} \varphi(y)$.  Then, we choose $b_{*} = \rho  - \beta r$. We note that $\nabla \varphi(x) \cdot \hat{x} \neq 0$. To see this, the supporting tangent plane of the epigraph of $\varphi(x)$ at $y$ is of the form $T(x) = a x + b$, then we note that $T(0) = b \leq 0$, since $T(x) \leq \varphi(x)$. But, if $a\cdot y = 0$, then $0<\varphi(y) = T(y) = b \leq 0$, which is a contradiction. Then, the line $\alpha r + b_{*}$ is below any other tangent line provided that the point of intersection at the origin is lower, which it is, and the slope is lower, which it is. Therefore, we have that $\alpha \vert x \vert + b_{*} \leq \varphi(x)$.

%Part E) Now, we do not assume smoothness. However, we assume that $\varphi(x) \neq \infty$. For this reason, the subdifferential of $\varphi(x)$ at a point $x$ is a compact set. We can repeat the argument above, but this time use subdifferentials instead.

%Part F) Suppose $\varphi(x) = \infty$ for some $x \in \mathbb{R}^n$, and assume $\varphi(x)$ achieves its finite minimum on a compact set. Then, there exists a convex function $\tilde{\varphi}(x)$ such that $\min_{x} \tilde{\varphi(x)} \neq -\infty$, the minimum of $\tilde{\varphi}$ occurs on a compact set, $\tilde{\varphi}(x) \leq \varphi(x)$, and $\tilde{\varphi}(x)$ is finite everywhere. Let $M>0$ and define the sublevel set $S_M := \{x \in \mathbb{R}^n : \varphi(x) \leq M \}$. This is a convex set that is compact, because the argument in part B) shows that $\varphi(x)$ cannot be a constant along any ray. Define the cone $C_{M} = \{ a(x,M) \in \mathbb{R}^{n+1} : a\geq 0, x \in S_{M} \}$. Then, define $\text{epi}(\varphi(x))_{M} = \{ (x, y) \in \mathbb{R}^{n+1} : x \in S_{M}, y \geq \varphi(x) \}$. Then, finally, define $G_{M} :=\{ \text{conv}(C_N \cup \text{epi}(\varphi(x))_{M}) \}$. We note $G_M$ is a convex set. It corresponds to the epigraph of a convex function $\tilde{\varphi}(x)$ that satisfies $\min_{x} \tilde{\varphi(x)} \neq -\infty$, the minimum of $\tilde{\varphi}$ occurs on a compact set, $\tilde{\varphi}(x) \leq \varphi(x)$ and $\tilde{\varphi}(x)$ is finite everywhere.

%Part D) Therefore, we conclude that $\varphi(x)$ must be proper and have a minimum in a compact subset $M \subset \mathbb{R}^n$. By an affine coordinate transformation we assume $\varphi(0) = 0$. Let $x_0$ be a point such that $\vert x_0 \vert >m$. For any point $y \in \vert x_0 \vert \mathbb{S}^{n-1}$, we find that the supporting tangent plane of the epigraph of $\varphi(x)$ at $y$ is of the form $T(x) = a x + b$, where $ay \neq 0$. To see this, we note that $T(0) = b \leq 0$, since $T(x) \leq \varphi(x)$. But, if $a\cdot y = 0$, then $0<\varphi(y) = T(y) = b \leq 0$, which is a contradiction. Since $\varphi(x)$ can take the value $+\infty$, at some points, the subgradient does not exist. At those points the subgradient is empty. Let $a^{*}$ denote that value $a^{*} = \inf_{y \in \vert x_0 \vert \mathbb{S}^{n-1}} \vert \partial \varphi(y)\vert$. Note that, since we are taking the minimum over a compact set of numbers which are greater than zero, we conclude $a^{*}>0$. Then, define $b^{*}$ via $b^{*} = -a^{*} \vert x_0 \vert$.

%Notice that since $\varphi(x)$ must satisfy $\int_{\mathbb{R}^n} e^{-\varphi(x)}dx = 1$, we cannot have $\text{epi}(\varphi) \subset \mathbb{R}^{n+1}$ have any recession directions in the closed half plane: $\left\{ (x_1, \dots, x_d, x_{d+1}) \vert x_{d+1} \leq 0 \right\}$. For if this were true, then the integral $\int_{\mathbb{R}^n} e^{-\varphi(x)}dx$ would necessarily be infinite. Therefore, the recession directions must be in the open upper half plane: $\left\{ (x_1, \dots, x_d, x_{d+1}) \vert x_{d+1} > 0 \right\}$. Since $\varphi(x)$ is bounded from below and since all the recession directions are in the open upper half plane let $x^{*} \in \mathbb{R}$ denote a point at which its minimum is achieved. Via an affine transformation, we move coordinates, so that $(x^{*}, \varphi(x^{*})$ is now at the origin in $\mathbb{R}^{n+1}$. Then, we define $K := \left\{\lambda \text{epi}(\varphi) \vert \lambda \geq 0, x\in \text{epi}(\varphi)  \right\}$. This is the smallest convex cone containing $\text{epi}(\varphi)$. Since $\text{epi}(\varphi)$ only has recession directions in the open upper half plane, the set $K$ satisfies $K \subset \text{epi}(g)$, for $g(x) = a\left\vert x \right\vert$, where $a>0$.

Part E) This then shows that there exist $a>0$ and $b \in \mathbb{R}$ such that $\varphi(x) \geq a\left\vert x \right\vert + b$. Thus, for log-concave distributions $\nu$, we can use Lemma~\ref{lemma:firstbound}. First, we compute:
\begin{align}\label{eq:mbn1}
1 - \nu(B_{R}(0)) &= \int_{\mathbb{R}^{n-1} \setminus B_{R}(0)} e^{-\varphi(x)}dx \\
&\leq \int_{\mathbb{R}^{n} \setminus B_{R}(0)} e^{-a \vert x \vert - b}dx \\
&\leq \vert \mathbb{S}^{n-1} \vert e^{-b} \int_{R}^{\infty} r^{n-1}e^{-a r}dr \\
&= \vert \mathbb{S}^{n-1} \vert e^{-b-aR} R^{n-1} \sum_{k=1}^{n} \frac{R^{1-k}}{a^k} \frac{(n-1)!}{(n-k)!}.
\end{align}
Therefore, for $R \geq 1$, we have:
\begin{align}
1 - \nu(B_{R}(0)) \leq \vert \mathbb{S}^{n-1} \vert e^{-b-aR} R^{n-1} \sum_{k=1}^{n} \frac{(n-1)!}{a^k(n-k)!}.
\end{align}
Likewise, we obtain the estimate, for $R \geq 1$:
\begin{equation}\label{eq:mbn2}
\int_{\mathbb{R}^n \setminus B_{R}(0)} \vert x \vert e^{-\varphi(x)}dx \leq \vert \mathbb{S}^{n-1} \vert e^{-b - aR} R^n \sum_{k=1}^{n+1} \frac{n!}{a^k (n - k+1)!}.
\end{equation}
Therefore, using Inequality~\eqref{eq:generalbound}, we obtain, for $R \geq 1$, that:
\begin{align}\label{eq:fastrate}
W_1(\nu, \nu_R) &\leq \left( \sum_{k=1}^{n+1} \frac{n!}{a^k(n - k + 1)!} + \sum_{k=1}^{n} \frac{(n-1)!}{a^k(n - k)!} \right) \vert \mathbb{S}^{n-1} \vert R^n e^{-b - aR} \\
&= C(a, b, n) R^n e^{-aR}.
\end{align}
\end{proof}

\end{document}
